% Páginas de entrada comunes a los libros de Razonamiento lógico (1 y 2)
\cleardoublepage
\thispagestyle{vacia}
{\sffamily\fontsize{26}{30}\selectfont\firalight Cómo se usa este libro\par}\vspace{2pt}
{\color{zolive}\rule{18mm}{1.6pt}}\par\vspace{8pt}

Este libro y la aplicación \marca\ son dos mitades de un mismo curso. En la app cada lección es una explicación que se
construye con el dedo, paso a paso; el libro sigue \textbf{la misma ruta} y va más despacio y más hondo: explica \emph{por qué}
funciona cada idea, resuelve ejemplos completos y te deja papel y lápiz para practicar como en el examen.

\vspace{4pt}
\begin{tabularx}{\linewidth}{@{}>{\leavevmode\sffamily\small\bfseries\color{zolive}}p{32mm}X@{}}
La ruta & Abre la lección: los pasos que vas a recorrer, los mismos de la app. El último es siempre un problema.\\[4pt]
La analogía & Una comparación con algo cotidiano para entender la idea antes de cualquier fórmula.\\[4pt]
Cuadro por cuadro & Cada idea se construye en dibujos numerados, como una historieta: se mira el cuadro, se lee su texto y
se sigue. La fórmula aparece cuando el dibujo ya la mostró.\\[4pt]
Tu turno & Enseguida, un ejercicio con el mismo dibujo: donde la app pide arrastrar, aquí se marca con lápiz. Compara con
\emph{Soluciones dibujadas}, al final.\\[4pt]
Ejemplo resuelto & Un problema completo, con un dibujo para cada paso de la solución.\\[4pt]
La regla & El recuadro con lo que hay que llevarse. Si solo tienes cinco minutos antes del examen, lee los recuadros.\\[4pt]
Ejemplos & Problemas resueltos. En cada uno hay una línea de tijera: tapa desde ahí, resuélvelo tú y luego compara.\\[4pt]
Un reto & Un problema como los del examen. Escribe tu respuesta y qué tan seguro estás; debajo, «Así se usó el concepto» muestra cómo se aplica.\\[4pt]
Arma la respuesta & No se pide un resultado: se pide construir algo que cumpla una condición. Las claves dicen el error que descubre.\\[4pt]
Lo que construiste & Las ideas de la lección, para marcar las que podrías explicar.\\[4pt]
Taller & Tres pisos: \textbf{Calienta}, \textbf{Entrena} y \textbf{Como en el examen} (opción múltiple A–D).\\[4pt]
Evaluación & Cada mundo termina con una prueba como la del examen, con tiempo.\\[4pt]
\end{tabularx}
{\small\color{zmute}Según la lección encontrarás también trampas (el error que más se repite), atajos y «Para ir más lejos».\par}

\clearpage
\thispagestyle{vacia}
{\sffamily\fontsize{26}{30}\selectfont\firalight La prueba de razonamiento\par}\vspace{2pt}
{\color{zolive}\rule{18mm}{1.6pt}}\par\vspace{8pt}

La prueba de admisión de la Universidad de Antioquia tiene dos componentes con el mismo peso: Razonamiento lógico y
Competencia lectora. En razonamiento hay alrededor de \textbf{40 preguntas} de opción múltiple con cuatro opciones,
\textbf{sin calculadora}. No pide fórmulas raras: pide \emph{ver} la estructura de un problema, organizar la información
y descartar con criterio. Por eso este libro insiste en tres hábitos.

\begin{enumerate}[leftmargin=1.6em,itemsep=4pt]
  \item \textbf{Dibujar antes de calcular.} Una recta, una tabla, un árbol o una barra partida convierten un enunciado
  enredado en algo que se ve. Casi todos los ejemplos del libro empiezan con un dibujo.
  \item \textbf{Calcular lo menos posible.} Simplificar antes de multiplicar, estimar antes de operar, probar las opciones
  cuando el problema lo permite. Una cuenta corta y segura vale más que una larga y frágil.
  \item \textbf{Leer la pregunta dos veces.} «¿Cuántos \emph{más}?», «\emph{a lo sumo}», «\emph{necesariamente}», «el
  \emph{resto}»: una palabra cambia la respuesta. Las opciones falsas del examen suelen ser la respuesta a otra
  pregunta muy parecida.
\end{enumerate}

\subsection*{Cómo repartir el tiempo}
Con unos 90 segundos por pregunta, conviene una primera vuelta rápida: responder lo que sale en menos de un minuto,
marcar con un círculo lo que parece largo y volver después. Ninguna pregunta vale más que otra; una pregunta difícil no
merece cinco minutos si hay tres fáciles esperando.

\subsection*{Cómo descartar}
\begin{itemize}
  \item \textbf{Por tamaño:} si el resultado debe estar entre 50 y 60, las opciones fuera de ese rango sobran.
  \item \textbf{Por paridad o última cifra:} un producto de un número par es par; $37 \times 43$ termina en 1.
  \item \textbf{Por unidades:} un área se mide en unidades cuadradas, un volumen en cúbicas.
  \item \textbf{Probando:} cuando piden un número que cumple condiciones, reemplazar las opciones suele ser más rápido que
  plantear la ecuación.
\end{itemize}

\subsection*{Símbolos que vas a encontrar}
\begin{multicols}{2}\small
\begin{tabular}{@{}>{$}l<{$}l@{}}
a \cdot b,\ ab & producto\\
\tfrac{a}{b} & cociente, fracción\\
a^n & potencia\\
\sqrt{a} & raíz cuadrada\\
|a| & valor absoluto\\
a \mid b & $a$ divide a $b$\\
a \approx b & aproximadamente igual\\
\end{tabular}

\begin{tabular}{@{}>{$}l<{$}l@{}}
a \neq b & distinto de\\
<,\ \le & menor, menor o igual\\
n! & factorial\\
\pi \approx 3.14 & número pi\\
\overline{ab} & número de cifras $a$, $b$\\
\angle ABC & ángulo de vértice $B$\\
\triangle ABC & triángulo\\
\end{tabular}
\end{multicols}
